% Part: many-valued-logic
% Chapter: syntax-and-semantics
% Section: introduction

\documentclass[../../../include/open-logic-section]{subfiles}

\begin{document}

\olfileid{mvl}{syn}{int}

\olsection{مقدمه}

در منطق کلاسیک با~!!{formula}هایی سروکار داریم که از
!!{propositional variable}ها و رابط‌های گزاره‌ای~$\lnot$، $\land$،
$\lor$، $\lif$ و~$\liff$ ساخته شده‌اند. معناشناسی منطق کلاسیک را با دو
ارزش صدق~$\True$ و~$\False$ تعریف می‌کنیم. !!{propositional variable}ها
را در~!!a{valuation} چون~$\pAssign{v}$ تفسیر می‌کنیم که این ارزش‌های
صدق~$\True$ و~$\False$ را به~!!{propositional variable}ها نسبت می‌دهد.
سپس هر~!!{valuation} برای هر~!!{formula} چون~$!A$، ارزش صدقی
$\pValue{v}(!A)$ تعیین می‌کند؛ و~!!^a{formula} در~!!a{valuation}
$\pAssign{v}$ ارضا می‌شود، یعنی~$\pSat{v}{!A}$، اگر و تنها اگر
$\pValue{v}(!A) = \True$.

منطق‌های چندارزشی تعمیم‌هایی از منطق کلاسیک دوارزشی‌اند که بیش از دو
ارزش صدق~$\True$ و~$\False$ را مجاز می‌دانند. بنابراین در منطق
چندارزشی، !!a{valuation} چون~$\pAssign{v}$ تابعی است که به هر
!!{propositional variable} چون~$p$ یکی از دامنه‌ای از ارزش‌های صدق ممکن
را نسبت می‌دهد. عموماً مجموعهٔ ارزش‌های صدق مجاز را~$V$ می‌نامیم. منطق
کلاسیک منطقی چندارزشی است که در آن~$V = \{\True, \False\}$؛ و ارزش
صدق~$\pValue{v}(!A)$ با جدول‌های صدق مشخصه و آشنای رابط‌ها محاسبه
می‌شود.

پس از افزودن ارزش‌های صدق دیگر، برای چگونگی محاسبهٔ
$\pValue{v}(!A)$ در مورد رابط‌هایی که «و»، «یا»، «نه» و «اگر---آنگاه»
می‌خوانیم بیش از یک انتخاب طبیعی داریم. پس یک منطق چندارزشی نه‌فقط با
مجموعهٔ ارزش‌های صدق، بلکه با \emph{توابع صدق}ی تعیین می‌شود که برای هر
رابط برمی‌گزینیم. بااین‌حال، پس از انتخاب این توابع برای منطق چندارزشی
$\Log L$، ارزش صدق~$\pValue{v}(!A)[\Log L]$، درست مانند منطق کلاسیک،
به‌طور یکتا با ارزش‌گذاری تعیین می‌شود. منطق‌های چندارزشی نیز مانند
منطق کلاسیک \emph{تابع‌ارزشی} هستند.

با داشتن این اجزای معناشناختی می‌توانیم همتاهای چندارزشی مفاهیم
معناشناختی همان‌گویی، استلزام و ارضاپذیری را تعریف کنیم. در منطق کلاسیک،
!!a{formula} هنگامی همان‌گویی است که برای هر~$\pAssign{v}$، ارزش صدقش
$\pValue{v}(!A) = \True$ باشد. در منطق چندارزشی باید این مفهوم را نیز
اندکی تعمیم دهیم. نخست، هیچ الزامی نیست که مجموعهٔ ارزش‌های صدق~$V$
شامل~$\True$ باشد. برای مثال، برخی منطق‌های چندارزشی مجموعه‌ای از اعداد،
مانند همهٔ اعداد گویا میان~$0$ و~$1$، را مجموعهٔ ارزش‌های صدق خود قرار
می‌دهند. در چنین حالتی معمولاً~$1$ نقش~$\True$ را دارد. در منطق‌های
دیگر، نه یک بلکه چند ارزش صدق این نقش را دارند. ازاین‌رو مقرر می‌کنیم
هر منطق چندارزشی مجموعه‌ای~$V^+$ از \emph{ارزش‌های ممتاز} داشته باشد.
سپس می‌توانیم بگوییم~!!a{formula} در~!!a{valuation} چون~$\pAssign{v}$
ارضا می‌شود، یعنی~$\pSat{v}{!A}[\Log L]$، اگر و تنها اگر
$\pValue{v}(!A)[\Log L] \in V^+$. !!^a{formula} چون~$!A$ همان‌گوییِ
منطق است، یعنی~$\Entails[\Log L] !A$، اگر و تنها اگر برای هر
$\pAssign{v}$، $\pValue{v}(!A)[\Log L] \in V^+$. سرانجام، می‌گوییم
$!A$ از مجموعه‌ای از~!!{formula}ها نتیجه می‌شود، یعنی
$\Gamma \Entails[\Log L] !A$، اگر هر~!!{valuation} که همهٔ
!!{formula}های~$\Gamma$ را ارضا می‌کند، $!A$ را نیز ارضا کند.

\end{document}
